% !PRINTED 556
\weberpage{556}

\weq{6}{\varphi(r) = (b_1 - b_0)\,r + (b_2 - b_0)\,r^2 + \cdots + (b_{n-2} - b_0)\,r^{n-2} - b_0 r^{n-1};}

oder da die Potenzen $r, r^2, \ldots r^{n-1}$, von der Reihenfolge abgesehen, mit $r, r_1, r_2 \ldots r_{n-2}$ übereinstimmen, in die Form

\weq{7}{\varphi(r) = a\,r + a_1 r_1 + a_2 r_2 + \cdots + a_{n-2} r_{n-2},}

worin die Coëfficienten $b$ und $a$ rationale Zahlen sind.

\gesperrt{I. Eine solche Function $\varphi(r)$ kann nur dann gleich Null sein, wenn alle ihre Coëfficienten $a, a_1 \ldots a_{n-2}$ verschwinden;}

denn der Ausdruck (5) zeigt, dass $\varphi(r)$ nicht anders verschwinden kann, als wenn die $b_0, b_1 \ldots b_{n-2}$ Null sind. Sind aber diese Null, so zeigt (6), dass auch die Coëfficienten $a$ alle Null sein müssen. Es kann also $\varphi(r)$ nur auf eine Weise in die Form (7) gebracht werden.

Treffen wir die Festsetzung, dass $a_h$ und $r_h$ sich nicht ändern sollen, wenn der Index $h$ um ein Vielfaches von $n - 1$ wächst, so können wir auch

\weq{8}{\varphi(r) = \sum^h a_h r_h}

setzen und darin $h$ ein volles Restsystem nach dem Modul $n - 1$ durchlaufen lassen.

Wir haben nun nachzuweisen, dass die Gruppe der Gleichung $X = 0$ keine andere ist, als die Periode der cyklischen Permutation
\[
\pi = (r, r_1, r_2 \ldots r_{n-2}).
\]

Diese Periode oder cyklische Gruppe bezeichnen wir mit

\weq{9}{C = 1, \pi, \pi^2 \ldots \pi^{n-2}.}

Die Gleichung $X = 0$ ist jedenfalls eine \gesperrt{Normalgleichung}, weil sie irreducibel ist, und weil alle ihre Wurzeln nach (3) rational durch eine unter ihnen ausdrückbar sind; sie ist also ihre eigene \gesperrt{Galois'sche} Resolvente (\S.\,145). Ihre Substitutionen sind

\weq{10}{(r, r), (r, r_1), (r, r_2) \ldots (r, r_{n-2}),}

und nach (3) ist
\[
(r, r_h) = (r_k, r_{k+h}).
\]

Daraus folgt, dass $(r, r_1)$ unter den Wurzeln (4) die cyklische Permutation $\pi$ hervorruft, und dass also die Substitutionsgruppe
%CONT>>
